\documentclass[10pt]{article}
\usepackage[left=1in,right=1in,top=3cm,bottom=0.35in,%
includeheadfoot,headheight=99pt,headsep=0.2in]{geometry}
\usepackage{times,graphicx,cite,fancyhdr,amsmath,amssymb,captdef}
\usepackage{marvosym}
\renewcommand{\headrulewidth}{0pt}
%----------------  page style -------------------
\pagestyle{fancy}
\fancyhead{}    %clear the head
\fancyhead[L]{
  \begin{minipage}[t]{1.8in}
    \hspace{0.6in}
    \includegraphics[height=0.6in]{anl.png}
  \end{minipage}
}
\fancyhead[C]{}
\fancyhead[R]{%
  \begin{minipage}[t]{1.7in}
    \begin{flushleft}
      {\scriptsize
        \vspace{-1.1in}
        {\bf Qingteng Zhang, Ph.D., Physicist}\\
        \vspace{0.08in}
        {\bf X-ray Science Division}\\
        Argonne National Laboratory \\
        9700 South Cass Avenue\\
        Argonne, IL 60439 \\
        \vspace{0.08in}
        \Telefon: 630-252-6366 \\
        }
    \end{flushleft}
  \end{minipage}
}
\fancyfoot{}    %clear the foot
\fancyfoot[C]{
  \footnotesize{A U.S. Department of Energy laboratory managed by UChicago
    Argonne, LLC}
}
% ==============================================================================
\setlength{\parskip}{2ex}
\setlength{\parindent}{0pt}

\usepackage[hidelinks]{hyperref}
\usepackage{enumitem}
\setlength{\parskip}{0.4em}
\setlength{\parindent}{0pt}
\raggedbottom
\sloppy   % long unbreakable e-mail addresses have no hyphenation points

\begin{document}

\today

Editors\\
\textit{Proceedings of the National Academy of Sciences}

Dear Editors:

Please consider the enclosed Research Report as a Direct Submission to
\textit{PNAS}: ``Dynamic Arrest Follows Assembly on a Temperature-Scaled
Common Clock in an Elastin-like Polypeptide.''

\textbf{Corresponding authors.} Qingteng Zhang
\mbox{(\href{mailto:qzhang234@anl.gov}{qzhang234@anl.gov})}, Arvind Ramanathan
\mbox{(\href{mailto:ramanathana@anl.gov}{ramanathana@anl.gov})}, and Gyorgy Babnigg
\mbox{(\href{mailto:gbabnigg@anl.gov}{gbabnigg@anl.gov})}.

\textbf{Classification.} Physical Sciences: Applied Physical Sciences.
Secondary: Biophysics and Computational Biology.

\textbf{Significance.} Whether a soft or biological material has stopped
flowing is usually inferred from its structure: a scattering profile that has
grown enough is read as a gel. We measure structure and motion on the same
detector frames and find that they run on one clock. In isothermal
experiments on the sequence-defined elastin-like polypeptide (VPAVG)$_{30}$ at
four temperatures between 28.00 and 28.75~$^\circ$C, every sample assembles,
but only those above 28.00~$^\circ$C arrest within the experiment. The
low-$Q$ SAXS intensity and the arrested fraction collapse with the same
scaling temperature, and arrest follows assembly: every sample that arrests
is half arrested at nearly the same intensity, and most of the
arrest happens after most of the low-$Q$ structure has formed. Arrest reaches
the largest length scales first, and once under way it advances by converting
material out of the mobile population rather than by slowing it: the
fast-decay amplitude falls from 0.80 to 0.07 at 175~nm within an hour while
the $1/e$ rate of the fast decay levels off.

\textbf{Novelty.} One rescaling of the waiting time collapses all four
isothermal experiments onto a single master curve: arrest speeds up by a factor
of e for every $T_s=0.24\pm0.02$~K of heating, and the low-$Q$ SAXS superposes with the
same $T_s$. That is a stronger statement than a rate: most of what
distinguishes the four experiments is their pace, and temperature sets it. A scaling temperature of a few tenths of a
kelvin is far too sharp for a single contact, which places the rate-limiting
step at the concerted formation of a junction in which many contacts close
together.

That scaling framework is our own. One of us introduced it for thermoreversible
colloidal gels [Q. Zhang \textit{et al.}, \textit{Phys. Rev. Lett.}
\textbf{119}, 178006 (2017)], where it returns $T_s=0.33$ and 0.45~K for two
particle sizes [D. Bahadur \textit{et al.}, \textit{J. Chem. Phys.}
\textbf{151}, 104902 (2019)]. The present work carries it from colloids to a
protein polymer and finds the same cooperative regime. We cite the agreement
with the nanoparticle-gel phenomenology rather than claiming it as new; what is
new is that structure and dynamics share one clock in a protein network, and
that one rescaling governs every temperature we can measure.

\textbf{Why it needed this measurement.} The fast relaxation is
sub-millisecond at these length scales and falls at or below the shortest
delay a conventional kilohertz acquisition reaches. Framing X-ray speckle at
50~kHz with the XSPA-500k at APS 8-ID-I resolves the fast decay itself: its
rate, its diffusive $Q$ dependence and its broadening as the sample arrests.
At conventional speeds only the loss of contrast would remain.

\textbf{Limits, stated plainly.} The 1.6~s maximum delay bounds a fraction
arrested on that scale rather than an infinite-time nonergodicity factor. The
scaling rests on three arresting temperatures spanning half a degree, and the
collapse is approximate rather than exact.

\textbf{Companion manuscript.} A separate manuscript, in preparation from the
same construct, reports the structural side: a thermally reversible transition
to a hierarchical assembly carrying $\beta$-sheet-like order, and the
dependence of that reversibility on chain length. The two were originally one
manuscript and were separated because the dynamical result stands on its own and
the structural result belongs with the chain-length series. There is no
overlap of figures or data between them. We are happy to supply the companion
manuscript on request.

\textbf{Data and code.} The reduced data and all analysis and figure-generation
code are public at
\href{https://github.com/qzhang234/PubCode-2021-IDPP}{qzhang234/PubCode-2021-IDPP};
every data-derived figure regenerates from that repository alone.

The manuscript is original, is not under consideration elsewhere, and has been
approved by all authors. The authors declare no competing interest;
AI-tool use is disclosed in the Acknowledgments. Suggested Editorial Board
members, NAS members and reviewers are listed overleaf.

Thank you for your consideration.

Sincerely,

\vspace{0.1in}
\nopagebreak
Qingteng Zhang, Ph.D., on behalf of all authors

\newpage
\section*{Suggested Editors and Reviewers}
\small
PNAS asks for a minimum of three Editorial Board members, three NAS members and
three qualified reviewers. \textbf{TO BE COMPLETED BEFORE SUBMISSION.} The
researchers below have relevant expertise in elastin-like polypeptides,
protein and peptide assembly, small-angle X-ray scattering, and X-ray photon
correlation spectroscopy, and were carried over from the earlier cover letter;
they are reviewer suggestions, and the Editorial Board and NAS member lists
still have to be chosen.

\textbf{Bradley D. Olsen}\\
Alexander and I. Michael Kasser (1960) Professor of Chemical Engineering, Massachusetts Institute of Technology\\
\href{mailto:bdolsen@mit.edu}{bdolsen@mit.edu}\\
Expertise: elastin-like polypeptide gels, associative protein polymers, arrested phase separation, and neutron spin-echo measurement of relaxation in protein hydrogels.

\textbf{Darrin J. Pochan}\\
Professor of Materials Science and Engineering, University of Delaware\\
\href{mailto:pochan@udel.edu}{pochan@udel.edu}\\
Expertise: peptide and protein self-assembly, $\beta$-sheet hydrogel networks, and small-angle scattering and rheology of soft biomaterials.

\textbf{Andrei Fluerasu}\\
Lead Beamline Scientist, Coherent Hard X-ray Scattering, National Synchrotron Light Source II, Brookhaven National Laboratory\\
\href{mailto:fluerasu@bnl.gov}{fluerasu@bnl.gov}\\
Expertise: X-ray photon correlation spectroscopy and the microscopic dynamics of gels and arrested soft matter.

\textbf{Theyencheri Narayanan}\\
Senior Scientist, European Synchrotron Radiation Facility, Grenoble\\
\href{mailto:narayan@esrf.fr}{narayan@esrf.fr}\\
Expertise: small-angle X-ray scattering and X-ray photon correlation spectroscopy of colloids and soft matter.

\end{document}
